% Source: https://artofproblemsolving.com/community/c3h2509055_hardest_math_problem_ever
% Stable identity: aops-topic-2509055-q01
Let $S(n)$ be the sum of the factorials of the digits of $n$.

\textbf{Bound.} If $n$ has $d$ digits, then $n\ge10^{d-1}$ and $S(n)\le d\cdot9!=362880d$. For $d=8$, $8\cdot362880=2903040<10^7$, and when $d$ increases by $1$ the bound $362880d$ grows by $362880$ while $10^{d-1}$ is multiplied by $10$. So $S(n)<n$ whenever $d\ge8$: a solution has at most $7$ digits, and $n\le7\cdot9!=2540160$.

\textbf{One and two digits.} $n=n!$ gives $n=1$ or $n=2$. For two digits, a digit $\ge5$ gives $S(n)\ge120>99$; so all digits are at most $4$ and $n=10a+b\le S(n)\le48$, i.e. $a\le4$. Then $a!+b!=10a+b$ means $b!=b+9$, $b+18$, $b+24$ or $b+16$ for $a=1,2,3,4$, and no $b\le4$ satisfies any of these.

\textbf{Three digits.} No digit is $7$ or more ($7!>999$). A digit $6$ would give $n=S(n)\ge720$, so the first digit would be at least $7$: impossible. If all digits are at most $4$, then $S(n)\le72<100$. So some digit is $5$, and $122\le S(n)\le360$, so the first digit $a$ is $1$, $2$ or $3$. If both other digits are $5$, then $S(n)=a!+240\in\{241,242,246\}$ while $n\in\{155,255,355\}$: no. Otherwise exactly one other digit is $5$ and the remaining digit $e$ is at most $4$, so $S(n)\le6+120+24=150$; then $a=1$ and $S(n)=121+e!\le145$, which rules out $n=15e\ge150$. So $n=1e5=105+10e$ with $105+10e=121+e!$, i.e. $e!=10e-16$, whose only solution with $e\le4$ is $e=4$: $n=145=1!+4!+5!$.

\textbf{Four to seven digits.} This is a finite check of the numbers from $1000$ to $2540160$ (it can be organised by the multiset of digits, since $S(n)$ depends only on the digits and not on their order); carried out by computer, it finds exactly one more solution:
\[40585=4!+0!+5!+8!+5!=24+1+120+40320+120.\]

So the numbers are $\boxed{1,\ 2,\ 145,\ 40585}$.
